\documentclass[conference]{IEEEtran}
\IEEEoverridecommandlockouts
% The preceding line is only needed to identify funding in the first footnote. If that is unneeded, please comment it out.
\usepackage{cite}
\usepackage{amsmath,amssymb,amsfonts}
\usepackage{algorithmic}
\usepackage{graphicx}
\usepackage{textcomp}
\usepackage{xcolor}
\usepackage[hidelinks]{hyperref}
\def\BibTeX{{\rm B\kern-.05em{\sc i\kern-.025em b}\kern-.08em
    T\kern-.1667em\lower.7ex\hbox{E}\kern-.125emX}}
\begin{document}

\title{Conference Paper Title}

\author{
\IEEEauthorblockN{Deniz Altunkapan}
\IEEEauthorblockA{\textit{University of Stuttgart}\\
Stuttgart, Germany \\
deniz.altunkapan@outlook.com}
\and
\IEEEauthorblockN{Maximilian Peresunchak}
\IEEEauthorblockA{\textit{University of Stuttgart}\\
Stuttgart, Germany \\
max.peresunchak@gmail.com}
\and
\IEEEauthorblockN{Viorel Tsigos}
\IEEEauthorblockA{\textit{University of Stuttgart}\\
Stuttgart, Germany \\
viorel.tsigos@gmail.com}
\and
\IEEEauthorblockN{Nico Reng}
\IEEEauthorblockA{\textit{University of Stuttgart}\\
Stuttgart, Germany \\
n.reng@gmail.com}
\and
\IEEEauthorblockN{Silas Aidam}
\IEEEauthorblockA{\textit{University of Stuttgart}\\
Stuttgart, Germany \\
silas@aidam.org}
\and
\IEEEauthorblockN{Michael Schramm}
\IEEEauthorblockA{\textit{University of Stuttgart}\\
Stuttgart, Germany \\
michael.schramm@ipvs.uni-stuttgart.de}
\and
\IEEEauthorblockN{Christian Becker}
\IEEEauthorblockA{\textit{University of Stuttgart}\\
Stuttgart, Germany \\
christian.becker@ipvs.uni-stuttgart.de}
}

\maketitle

\begin{abstract}
Federated Learning (FL) trains a shared global model by iterating between local client-side optimization and server-side aggregation of model updates, while raw data remains on client devices or at participating institutions. This training paradigm is attractive for privacy-sensitive and cross-organization settings, but practical deployments are often geo-distributed and heterogeneous. As a result, communication delays, slow-client bottlenecks, and unstable runtime conditions can significantly affect convergence and total training time. Multi-server approaches, including hierarchical FL and asynchronous designs such as Spyker, are proposed to reduce these limitations, yet reproducible and topology-aware evaluation tooling remains limited in practice.

This paper presents a graph-configurable simulation framework for geo-distributed multi-server FL experimentation. The framework is built on top of Flower and extends it with support for hierarchical multi-tier aggregation and asynchronous multi-server coordination (including Spyker-style execution). It integrates a web-based topology editor, run lifecycle management, structured runtime logging, and observability hooks for metric inspection. The system is designed for controlled experimentation through explicit configuration of topology, runtime parameters, and seeds, while enforcing architecture-aware validation at both UI and backend levels.

Beyond implementation, the work positions the framework as an evaluation vehicle for comparing established FL simulation ecosystems with respect to multi-server support, topology expressiveness, scalability, observability, and extensibility. The resulting artifact provides a practical basis for reproducible benchmarking and method comparison in geo-distributed FL research settings.
\end{abstract}

\begin{IEEEkeywords}
federated learning, geo-distributed systems, multi-server simulation
\end{IEEEkeywords}

\section{Introduction}
Federated Learning (FL) has matured from a conceptual privacy-preserving training paradigm into a practical systems problem. The core optimization loop is well known: clients train locally on private data, send model updates, and a coordinator aggregates them into a new global model. However, once deployments become geo-distributed, FL behavior is no longer driven by learning dynamics alone. Network heterogeneity, regional latency asymmetries, and runtime resource contention become first-order factors for convergence speed and stability \cite{kairouz2021}.

This systems perspective has motivated multi-server approaches. Hierarchical FL introduces intermediate aggregation tiers to shorten communication paths and reduce synchronization pressure \cite{abad2020}. Asynchronous multi-server designs, including Spyker, relax strict round barriers to better tolerate slow clients and heterogeneous client completion times \cite{zuo2024}. These methods broaden the design space, but they also raise the bar for experimental infrastructure: researchers need simulation environments that can represent topology explicitly, execute architecture-specific control logic, and preserve traceability across runs.

Existing frameworks provide valuable capabilities, but differ in their emphasis. Flower and FLSim offer flexible FL experimentation interfaces \cite{flower,flsim}; pfl-research and Pollen emphasize simulation throughput and large-scale studies \cite{granqvist2024,sani2023}. For geo-distributed multi-server questions, the practical challenge is to combine three requirements in one workflow: topology-aware configuration, robust and reproducible run control, and system-level observability during execution.

This paper presents a framework designed around exactly these requirements. It is implemented as a Flower-based runtime extended with topology-aware orchestration components. The framework models deployments as editable graphs, executes hierarchical and Spyker-style asynchronous multi-server simulations through a unified run interface, and tracks each run via explicit lifecycle states (queued, running, completed, failed). The implementation couples frontend preflight checks with backend-authoritative validation, enforces consistent run identifiers and metadata propagation, and integrates structured logging and monitoring hooks for online inspection and post-hoc analysis.

Beyond presenting the framework itself, the paper positions it as an evaluation vehicle for systematic framework-level comparison under geo-distributed assumptions. The goal is not to claim algorithmic novelty, but to provide an experimentally reliable basis for comparing architectures and tools under controlled, reproducible conditions.

The main contributions of this work are:
\begin{itemize}
\item A graph-based, topology-aware simulation environment for geo-distributed multi-server FL experiments.
\item Support for hierarchical and asynchronous multi-server execution modes within a unified configuration and run interface.
\item A reproducibility- and observability-oriented workflow including run lifecycle tracking, validation, and structured logging integration for dashboard-based analysis.
\item A basis for systematic framework comparison using consistent criteria for multi-server capability, topology modeling, scalability, observability, and extensibility.
\end{itemize}

The remainder of the paper details the system design, implementation decisions, evaluation protocol, comparative analysis setup, and limitations.

\section{Ease of Use}

\subsection{Maintaining the Integrity of the Specifications}

The IEEEtran class file is used to format your paper and style the text. All margins, 
column widths, line spaces, and text fonts are prescribed; please do not 
alter them. You may note peculiarities. For example, the head margin
measures proportionately more than is customary. This measurement 
and others are deliberate, using specifications that anticipate your paper 
as one part of the entire proceedings, and not as an independent document. 
Please do not revise any of the current designations.

\section{Prepare Your Paper Before Styling}
Before you begin to format your paper, first write and save the content as a 
separate text file. Complete all content and organizational editing before 
formatting. Please note sections \ref{AA}--\ref{SCM} below for more information on 
proofreading, spelling and grammar.

Keep your text and graphic files separate until after the text has been 
formatted and styled. Do not number text heads---{\LaTeX} will do that 
for you.

\subsection{Abbreviations and Acronyms}\label{AA}
Define abbreviations and acronyms the first time they are used in the text, 
even after they have been defined in the abstract. Abbreviations such as 
IEEE, SI, MKS, CGS, ac, dc, and rms do not have to be defined. Do not use 
abbreviations in the title or heads unless they are unavoidable.

\subsection{Units}
\begin{itemize}
\item Use either SI (MKS) or CGS as primary units. (SI units are encouraged.) English units may be used as secondary units (in parentheses). An exception would be the use of English units as identifiers in trade, such as ``3.5-inch disk drive''.
\item Avoid combining SI and CGS units, such as current in amperes and magnetic field in oersteds. This often leads to confusion because equations do not balance dimensionally. If you must use mixed units, clearly state the units for each quantity that you use in an equation.
\item Do not mix complete spellings and abbreviations of units: ``Wb/m\textsuperscript{2}'' or ``webers per square meter'', not ``webers/m\textsuperscript{2}''. Spell out units when they appear in text: ``. . . a few henries'', not ``. . . a few H''.
\item Use a zero before decimal points: ``0.25'', not ``.25''. Use ``cm\textsuperscript{3}'', not ``cc''.)
\end{itemize}

\subsection{Equations}
Number equations consecutively. To make your 
equations more compact, you may use the solidus (~/~), the exp function, or 
appropriate exponents. Italicize Roman symbols for quantities and variables, 
but not Greek symbols. Use a long dash rather than a hyphen for a minus 
sign. Punctuate equations with commas or periods when they are part of a 
sentence, as in:
\begin{equation}
a+b=\gamma\label{eq}
\end{equation}

Be sure that the 
symbols in your equation have been defined before or immediately following 
the equation. Use ``\eqref{eq}'', not ``Eq.~\eqref{eq}'' or ``equation \eqref{eq}'', except at 
the beginning of a sentence: ``Equation \eqref{eq} is . . .''

\subsection{\LaTeX-Specific Advice}

Please use ``soft'' (e.g., \verb|\eqref{Eq}|) cross references instead
of ``hard'' references (e.g., \verb|(1)|). That will make it possible
to combine sections, add equations, or change the order of figures or
citations without having to go through the file line by line.

Please don't use the \verb|{eqnarray}| equation environment. Use
\verb|{align}| or \verb|{IEEEeqnarray}| instead. The \verb|{eqnarray}|
environment leaves unsightly spaces around relation symbols.

Please note that the \verb|{subequations}| environment in {\LaTeX}
will increment the main equation counter even when there are no
equation numbers displayed. If you forget that, you might write an
article in which the equation numbers skip from (17) to (20), causing
the copy editors to wonder if you've discovered a new method of
counting.

{\BibTeX} does not work by magic. It doesn't get the bibliographic
data from thin air but from .bib files. If you use {\BibTeX} to produce a
bibliography you must send the .bib files. 

{\LaTeX} can't read your mind. If you assign the same label to a
subsubsection and a table, you might find that Table I has been cross
referenced as Table IV-B3. 

{\LaTeX} does not have precognitive abilities. If you put a
\verb|\label| command before the command that updates the counter it's
supposed to be using, the label will pick up the last counter to be
cross referenced instead. In particular, a \verb|\label| command
should not go before the caption of a figure or a table.

Do not use \verb|\nonumber| inside the \verb|{array}| environment. It
will not stop equation numbers inside \verb|{array}| (there won't be
any anyway) and it might stop a wanted equation number in the
surrounding equation.

\subsection{Some Common Mistakes}\label{SCM}
\begin{itemize}
\item The word ``data'' is plural, not singular.
\item The subscript for the permeability of vacuum $\mu_{0}$, and other common scientific constants, is zero with subscript formatting, not a lowercase letter ``o''.
\item In American English, commas, semicolons, periods, question and exclamation marks are located within quotation marks only when a complete thought or name is cited, such as a title or full quotation. When quotation marks are used, instead of a bold or italic typeface, to highlight a word or phrase, punctuation should appear outside of the quotation marks. A parenthetical phrase or statement at the end of a sentence is punctuated outside of the closing parenthesis (like this). (A parenthetical sentence is punctuated within the parentheses.)
\item A graph within a graph is an ``inset'', not an ``insert''. The word alternatively is preferred to the word ``alternately'' (unless you really mean something that alternates).
\item Do not use the word ``essentially'' to mean ``approximately'' or ``effectively''.
\item In your paper title, if the words ``that uses'' can accurately replace the word ``using'', capitalize the ``u''; if not, keep using lower-cased.
\item Be aware of the different meanings of the homophones ``affect'' and ``effect'', ``complement'' and ``compliment'', ``discreet'' and ``discrete'', ``principal'' and ``principle''.
\item Do not confuse ``imply'' and ``infer''.
\item The prefix ``non'' is not a word; it should be joined to the word it modifies, usually without a hyphen.
\item There is no period after the ``et'' in the Latin abbreviation ``et al.''.
\item The abbreviation ``i.e.'' means ``that is'', and the abbreviation ``e.g.'' means ``for example''.
\end{itemize}
An excellent style manual for science writers is \cite{b7}.

\subsection{Authors and Affiliations}
\textbf{The class file is designed for, but not limited to, six authors.} A 
minimum of one author is required for all conference articles. Author names 
should be listed starting from left to right and then moving down to the 
next line. This is the author sequence that will be used in future citations 
and by indexing services. Names should not be listed in columns nor group by 
affiliation. Please keep your affiliations as succinct as possible (for 
example, do not differentiate among departments of the same organization).

\subsection{Identify the Headings}
Headings, or heads, are organizational devices that guide the reader through 
your paper. There are two types: component heads and text heads.

Component heads identify the different components of your paper and are not 
topically subordinate to each other. Examples include Acknowledgments and 
References and, for these, the correct style to use is ``Heading 5''. Use 
``figure caption'' for your Figure captions, and ``table head'' for your 
table title. Run-in heads, such as ``Abstract'', will require you to apply a 
style (in this case, italic) in addition to the style provided by the drop 
down menu to differentiate the head from the text.

Text heads organize the topics on a relational, hierarchical basis. For 
example, the paper title is the primary text head because all subsequent 
material relates and elaborates on this one topic. If there are two or more 
sub-topics, the next level head (uppercase Roman numerals) should be used 
and, conversely, if there are not at least two sub-topics, then no subheads 
should be introduced.

\subsection{Figures and Tables}
\paragraph{Positioning Figures and Tables} Place figures and tables at the top and 
bottom of columns. Avoid placing them in the middle of columns. Large 
figures and tables may span across both columns. Figure captions should be 
below the figures; table heads should appear above the tables. Insert 
figures and tables after they are cited in the text. Use the abbreviation 
``Fig.~\ref{fig}'', even at the beginning of a sentence.

\begin{table}[htbp]
\caption{Table Type Styles}
\begin{center}
\begin{tabular}{|c|c|c|c|}
\hline
\textbf{Table}&\multicolumn{3}{|c|}{\textbf{Table Column Head}} \\
\cline{2-4} 
\textbf{Head} & \textbf{\textit{Table column subhead}}& \textbf{\textit{Subhead}}& \textbf{\textit{Subhead}} \\
\hline
copy& More table copy$^{\mathrm{a}}$& &  \\
\hline
\multicolumn{4}{l}{$^{\mathrm{a}}$Sample of a Table footnote.}
\end{tabular}
\label{tab1}
\end{center}
\end{table}

\begin{figure}[htbp]
\centerline{\fbox{\rule{0pt}{1.2in}\rule{0.9\columnwidth}{0pt}}}
\caption{Example of a figure caption.}
\label{fig}
\end{figure}

Figure Labels: Use 8 point Times New Roman for Figure labels. Use words 
rather than symbols or abbreviations when writing Figure axis labels to 
avoid confusing the reader. As an example, write the quantity 
``Magnetization'', or ``Magnetization, M'', not just ``M''. If including 
units in the label, present them within parentheses. Do not label axes only 
with units. In the example, write ``Magnetization (A/m)'' or ``Magnetization 
\{A[m(1)]\}'', not just ``A/m''. Do not label axes with a ratio of 
quantities and units. For example, write ``Temperature (K)'', not 
``Temperature/K''.

\section*{Acknowledgment}

The preferred spelling of the word ``acknowledgment'' in America is without 
an ``e'' after the ``g''. Avoid the stilted expression ``one of us (R. B. 
G.) thanks $\ldots$''. Instead, try ``R. B. G. thanks$\ldots$''. Put sponsor 
acknowledgments in the unnumbered footnote on the first page.

\section*{References}

Please number citations consecutively within brackets \cite{b1}. The 
sentence punctuation follows the bracket \cite{b2}. Refer simply to the reference 
number, as in \cite{b3}---do not use ``Ref. \cite{b3}'' or ``reference \cite{b3}'' except at 
the beginning of a sentence: ``Reference \cite{b3} was the first $\ldots$''

Number footnotes separately in superscripts. Place the actual footnote at 
the bottom of the column in which it was cited. Do not put footnotes in the 
abstract or reference list. Use letters for table footnotes.

Unless there are six authors or more give all authors' names; do not use 
``et al.''. Papers that have not been published, even if they have been 
submitted for publication, should be cited as ``unpublished'' \cite{b4}. Papers 
that have been accepted for publication should be cited as ``in press'' \cite{b5}. 
Capitalize only the first word in a paper title, except for proper nouns and 
element symbols.

For papers published in translation journals, please give the English 
citation first, followed by the original foreign-language citation \cite{b6}.

\begin{thebibliography}{00}
\bibitem{kairouz2021} P. Kairouz \textit{et al.}, ``Advances and open problems in federated learning,'' \textit{Foundations and Trends in Machine Learning}, vol. 14, no. 1--2, pp. 1--210, 2021.
\bibitem{abad2020} M. S. H. Abad \textit{et al.}, ``Hierarchical federated learning across heterogeneous cellular networks,'' in \textit{Proc. IEEE ICASSP}, 2020, pp. 8866--8870.
\bibitem{zuo2024} Y. Zuo \textit{et al.}, ``Spyker: Asynchronous multi-server federated learning for geo-distributed clients,'' in \textit{Proc. 25th Int. Middleware Conf.}, 2024, pp. 367--378.
\bibitem{flower} Flower AI, ``Flower,'' https://flower.ai, accessed: 2026-04-02.
\bibitem{flsim} Meta Research, ``FLSim,'' https://github.com/facebookresearch/FLSim, accessed: 2026-04-02.
\bibitem{granqvist2024} F. Granqvist \textit{et al.}, ``pfl-research: Simulation framework for accelerating research in private federated learning,'' in \textit{Advances in Neural Information Processing Systems}, vol. 37, 2024.
\bibitem{sani2023} L. Sani \textit{et al.}, ``Pollen: High-throughput federated learning simulation via resource-aware client placement,'' \textit{arXiv preprint arXiv:2306.17453}, 2023.
\bibitem{b1} G. Eason, B. Noble, and I. N. Sneddon, ``On certain integrals of Lipschitz-Hankel type involving products of Bessel functions,'' Phil. Trans. Roy. Soc. London, vol. A247, pp. 529--551, April 1955.
\bibitem{b2} J. Clerk Maxwell, A Treatise on Electricity and Magnetism, 3rd ed., vol. 2. Oxford: Clarendon, 1892, pp.68--73.
\bibitem{b3} I. S. Jacobs and C. P. Bean, ``Fine particles, thin films and exchange anisotropy,'' in Magnetism, vol. III, G. T. Rado and H. Suhl, Eds. New York: Academic, 1963, pp. 271--350.
\bibitem{b4} K. Elissa, ``Title of paper if known,'' unpublished.
\bibitem{b5} R. Nicole, ``Title of paper with only first word capitalized,'' J. Name Stand. Abbrev., in press.
\bibitem{b6} Y. Yorozu, M. Hirano, K. Oka, and Y. Tagawa, ``Electron spectroscopy studies on magneto-optical media and plastic substrate interface,'' IEEE Transl. J. Magn. Japan, vol. 2, pp. 740--741, August 1987 [Digests 9th Annual Conf. Magnetics Japan, p. 301, 1982].
\bibitem{b7} M. Young, The Technical Writer's Handbook. Mill Valley, CA: University Science, 1989.
\end{thebibliography}
\vspace{12pt}
\color{red}
IEEE conference templates contain guidance text for composing and formatting conference papers. Please ensure that all template text is removed from your conference paper prior to submission to the conference. Failure to remove the template text from your paper may result in your paper not being published.

\end{document}
